\documentclass[11pt]{article}
\pdfinfoomitdate=1
\pdftrailerid{}
\pdfsuppressptexinfo=15
\usepackage[margin=1in]{geometry}
\usepackage{amsmath,amssymb,amsthm,mathtools,booktabs,longtable,microtype}
\usepackage[T1]{fontenc}
\usepackage{lmodern}
\usepackage{xcolor}
\usepackage[colorlinks=true,linkcolor=blue!55!black,citecolor=blue!55!black,urlcolor=blue!55!black]{hyperref}
\newtheorem{theorem}{Theorem}[section]
\newtheorem{lemma}[theorem]{Lemma}
\newtheorem{corollary}[theorem]{Corollary}
\newtheorem{definition}[theorem]{Definition}
\newtheorem{remark}[theorem]{Remark}
\newtheorem{firewall}[theorem]{Firewall}
\newcommand{\Fcal}{\mathcal F}
\newcommand{\Dcal}{\mathfrak D}
\newcommand{\CSHT}{\mathrm{CSHT}_{67}}
\newcommand{\RBLPTE}{\mathrm{RBLPTE}_{67}}
\newcommand{\SPCC}{\mathrm{SPCC}_{67}}
\title{Dickman Localization of the Factor--67 Critical Saddle:\\Power-Scale Positivity, Root Bellman Closure, and the Subpower Core}
\author{Agentic Polymath Project}
\date{August 18, 2026}
\begin{document}
\maketitle

\begin{abstract}
The factor--67 annular route has an exact native rough expansion
\[
 \Fcal(Y,z)=\sum_{P^-(m)\ge z}\mu(m)m^{-1/2}b(Y/m),
\]
where $b$ is the complete finite-$P_{61}$ annular $5{:}3$ scalar.  Previous
work proves the exact first-owner ledger, the compulsory $(R-A)F$ residual,
and a no-go theorem for every subcritical count-depth stopping line.  We prove
here that every state whose least allowed rough prime is a fixed positive power
of its endpoint is unconditionally positive.

More precisely, if $Y^\theta\le z\le Y$ and
$u=\log Y/\log z$, then, uniformly for fixed $\theta>0$,
\[
 \Fcal(Y,z)=a\sqrt Y\,\rho(u)+o_\theta(\sqrt Y),
\]
where $a>0$ is the finite-$P_{61}$ square-root coefficient and $\rho$ is the
Dickman function.  We derive $\rho$ directly as the alternating rough-product
simplex and prove its positivity from the integral identity
$u\rho(u)=\int_{u-1}^u\rho(v)\,dv$.  Consequently every fixed power-scale state
satisfies the state-wise root Bellman inequality, and admits a one-dimensional
one-use scalar transport.

This does not prove the stronger target-preserving critical-saddle transport,
and it does not apply at the root threshold $67$, where
$\log67/\log Y\to0$.  It localizes every possible obstruction to the subpower
least-prime, harmonic-saddle core.  The Riemann Hypothesis remains unproved.
\end{abstract}

\begin{center}
\fbox{\parbox{0.91\textwidth}{\textbf{Status.} The power-scale theorem is
unconditional.  Full $\CSHT$, root $\RBLPTE$, the subpower core, and RH remain
open.}}
\end{center}

\tableofcontents

\section{Frozen inputs and the exact question}

Let $b(Y)=F_{61}(Y)$ be the complete grouped finite-$P_{61}$ annular $5{:}3$
scalar.  The controlling finite theorem gives
\begin{equation}\label{eq:base}
 b(Y)=a\sqrt Y+c+O(Y^{-3/2}),
 \qquad
 a=12\prod_{p\le61}\left(1-{1\over p}\right)>0,
\end{equation}
and $b(Y)\ge0$ for every real $Y\ge1$.

For a real least-prime threshold $z\ge67$, define the literal completed native
state
\begin{equation}\label{eq:state}
 \Fcal(Y,z)=
 \sum_{\substack{m\text{ squarefree}\\P^-(m)\ge z}}
 {\mu(m)\over\sqrt m}b(Y/m),
\end{equation}
with zero extension below the support of $b$.  Every term in
\eqref{eq:state} carries its native coefficient, product activation and
cumulative parity.  No Euler product without the product boundary is inserted.

The root state has $z=67$.  The present paper studies the terminal regime
$z\ge Y^\theta$ for fixed $\theta>0$ and asks whether it can carry a negative
obstruction.

\section{The exact alternating simplex}

For $u\ge0$, put $I_0(u)=1$ and
\begin{equation}\label{eq:Ik}
 I_k(u)={1\over k!}
 \int_{\substack{s_1,\ldots,s_k\ge1\\s_1+\cdots+s_k\le u}}
 {ds_1\cdots ds_k\over s_1\cdots s_k}
 \qquad(k\ge1).
\end{equation}
The integral vanishes for $k>u$.  Define
\begin{equation}\label{eq:D}
 \Dcal(u)=\sum_{k=0}^{\lfloor u\rfloor}(-1)^kI_k(u).
\end{equation}

\begin{theorem}[Dickman simplex identity]\label{thm:dickman}
The function $\Dcal$ is continuous, equals one on $[0,1]$, and satisfies
\[
 u\Dcal'(u)=-\Dcal(u-1)
 \qquad(u>1)
\]
away from integer knots.  Moreover
\begin{equation}\label{eq:dickman-integral}
 \boxed{u\Dcal(u)=\int_{u-1}^{u}\Dcal(v)\,dv}
 \qquad(u\ge1),
\end{equation}
and hence
\begin{equation}\label{eq:dickman-positive}
 \boxed{\Dcal(u)>0\quad\text{for every finite }u.}
\end{equation}
Thus $\Dcal$ is the Dickman function $\rho$.
\end{theorem}

\begin{proof}
Differentiating $I_k$ through the moving simplex boundary and multiplying by
$u=s_1+\cdots+s_k$, symmetry gives
\[
 uI_k'(u)=I_{k-1}(u-1).
\]
Alternating summation proves the delay equation.  The derivative of the two
sides of \eqref{eq:dickman-integral} is the same, and they agree at $u=1$.
Positivity follows interval by interval from the positive integral average in
\eqref{eq:dickman-integral}.
\end{proof}

The theorem is source-relevant: $I_k(u)$ is the continuum limit of the total
harmonic weight of squarefree $k$-prime products with every prime at least $z$
and product at most $z^u$.

\section{Uniform lower-order control}

Fix $\theta\in(0,1)$.  Throughout this section
\[
 Y^\theta\le z\le Y,
 \qquad
 u={\log Y\over\log z}\in[1,1/\theta].
\]
Every active rough product has at most
$K_\theta=\lfloor1/\theta\rfloor$ prime factors.

For $k\ge1$, define
\[
 A_k(Y,z)=
 \sum_{\substack{z\le p_1<\cdots<p_k\\p_1\cdots p_k\le Y}}
 (p_1\cdots p_k)^{-1/2}.
\]

\begin{lemma}[Uniform half-weight bound]\label{lem:halfweight}
For every fixed $\theta>0$ and $1\le k\le K_\theta$,
\[
 A_k(Y,z)\ll_\theta{\sqrt Y\over\log Y}
\]
uniformly in $Y^\theta\le z\le Y$.
\end{lemma}

\begin{proof}
For $k=1$ this is partial summation from the prime number theorem.  For
$k\ge2$, selecting one prime and dropping the ordering constraint gives
\[
 kA_k(Y,z)
 \le\sum_{z\le p\le Y/z^{k-1}}p^{-1/2}A_{k-1}(Y/p,z).
\]
In this range $\log(Y/p)\ge(k-1)\theta\log Y$.  The induction hypothesis and
Mertens' bound for $\sum1/p$ give the result.
\end{proof}

It follows that
\begin{equation}\label{eq:lowerorder}
 \sum_m m^{-1/2}\ll_\theta{\sqrt Y\over\log Y},
 \qquad
 \sum_m m\ll_\theta{Y^2\over\log Y},
\end{equation}
for the nonunit active rough products.  The second estimate follows from
$m\le Y^{3/2}m^{-1/2}$.

\section{Prime-harmonic convergence}

Define the finite measure on $[1,1/\theta]$
\[
 \nu_{Y,z}=\sum_{z\le p\le Y}{1\over p}
 \delta_{\log p/\log z}.
\]
Mertens' theorem gives, uniformly for
$1\le\alpha\le\beta\le1/\theta$,
\begin{equation}\label{eq:measure}
 \nu_{Y,z}([\alpha,\beta])
 =\log(\beta/\alpha)+o_\theta(1).
\end{equation}
Thus $\nu_{Y,z}$ converges weakly and uniformly to $ds/s$.  The total harmonic
weight of repeated prime coordinates is
\[
 O\left(\sum_{p\ge z}p^{-2}\right)=o_\theta(1).
\]
Consequently, for each fixed $k$, the squarefree degree-$k$ harmonic sum with
product at most $Y$ tends uniformly to $I_k(u)$.

\section{Power-scale native positivity}

\begin{theorem}[Uniform Dickman asymptotic]\label{thm:main}
For every fixed $\theta\in(0,1)$, uniformly for $Y^\theta\le z\le Y$,
\begin{equation}\label{eq:main}
 \boxed{
 \Fcal(Y,z)=a\sqrt Y\,\rho\!\left({\log Y\over\log z}\right)
 +o_\theta(\sqrt Y).
 }
\end{equation}
In particular, for all sufficiently large $Y$,
\begin{equation}\label{eq:positive}
 \boxed{
 \Fcal(Y,z)\ge {a\over2}\rho(1/\theta)\sqrt Y>0
 }
\end{equation}
uniformly in the same range.
\end{theorem}

\begin{proof}
Use \eqref{eq:base} when $Y/m$ is beyond the fixed activation range.  The
square-root term contributes
\[
 a\sqrt Y
 \sum_{\substack{m\le Y\\m\text{ squarefree}\\P^-(m)\ge z}}{\mu(m)\over m}.
\]
The constant term, the compact boundary where $Y/m$ is bounded, and the
$O((Y/m)^{-3/2})$ remainder are all
$O_\theta(\sqrt Y/\log Y)$ by Lemma~\ref{lem:halfweight} and
\eqref{eq:lowerorder}.

Expanding the remaining harmonic sum by rough degree, using
\eqref{eq:measure}, and discarding repeated coordinates at $o_\theta(1)$ cost
gives exactly \eqref{eq:D}.  Theorem~\ref{thm:dickman} now yields
\eqref{eq:main}.  Since $\rho$ decreases on compact positive intervals,
$\rho(u)\ge\rho(1/\theta)>0$, proving \eqref{eq:positive}.
\end{proof}

\begin{corollary}[All fixed power thresholds are safe]\label{cor:alltheta}
The earlier condition $\theta>e^{-1}$ is unnecessary.  For every fixed
$\theta>0$, every state with least allowed rough prime at least $Y^\theta$ is
eventually positive.
\end{corollary}

\section{Bellman and transport consequences}

At any state, retain an exact largest-prime decomposition
\[
 U_{\rm full}=U_Z-\mathfrak I_Z-\mathfrak T_Z.
\]
The state-wise root Bellman inequality is
\[
 \mathfrak T_Z\le U_Z-\mathfrak I_Z.
\]

\begin{corollary}[Power-sector root Bellman closure]\label{cor:rblpte}
For every fixed $\theta>0$, the state-wise root Bellman inequality holds for
all sufficiently large states satisfying $z\ge Y^\theta$.
\end{corollary}

\begin{proof}
The Bellman inequality is algebraically equivalent to $U_{\rm full}\ge0$.
Apply Theorem~\ref{thm:main}.
\end{proof}

At scalar scope, write the completed source as nonnegative even scalar masses
$e_i$ and odd scalar masses $o_j$.  Equation \eqref{eq:positive} gives
$\sum e_i\ge\sum o_j$, so a greedy one-dimensional flow covers all odd scalar
mass while using every even scalar mass at most once.

\begin{firewall}[Scalar transport is not $\CSHT$]
The preceding flow does not preserve target, atom ratios, row coordinates or
activation germs.  Positive zero-hinge scalar slack does not imply target
capacity.  Therefore the theorem closes the minimal scalar state, not the
stronger all-coordinate Hall producer.
\end{firewall}

\section{The subpower critical core}

Suppose a sequence of states $(Y_n,z_n)$ violates state-wise scalar positivity.
If $\log z_n/\log Y_n$ were bounded below by one fixed positive $\theta$, the
sequence would contradict Theorem~\ref{thm:main}.  Hence necessarily
\begin{equation}\label{eq:subpower}
 \boxed{{\log z_n\over\log Y_n}\longrightarrow0.}
\end{equation}
The prior subcritical-depth no-go independently shows that a successful
stopping line cannot remain at depth
$L\log L=o(\log\log Y)$.  Thus every possible obstruction lies simultaneously
in the subpower least-prime regime and the rough harmonic saddle.

\begin{definition}[Subpower critical core, $\SPCC$]
$\SPCC$ is eventual nonnegativity of the literal native scalar on all state
sequences satisfying \eqref{eq:subpower} and lying in the harmonic-saddle depth
regime.
\end{definition}

The exact implication is
\[
 \SPCC\Longrightarrow\RBLPTE
 \Longrightarrow\text{eventual annular }5{:}3\text{ positivity}
 \Longrightarrow\mathrm{RH}.
\]
This paper does not prove $\SPCC$, root $\RBLPTE$, full $\CSHT$, or RH.  Its
unconditional contribution is that no fixed power-scale terminal state can
carry the obstruction.

\section{Replay and scientific boundary}

The deterministic replay checks exact residual and exposure identities,
computes the Dickman delay equation numerically for diagnostics, and evaluates
the literal native state at $Y=10^6$ for thresholds corresponding to
$\theta=1/2,2/5,1/3$.  All three states are positive and their rough harmonic
sums approach the Dickman values.  These computations are reconnaissance; the
proof is the preceding analytic argument.

Expected classification:
\[
 \texttt{PASS\_T97900\_DICKMAN\_POWER\_SECTOR\_LOCALIZATION}.
\]

\begin{center}
\begin{tabular}{ll}
\toprule
Dickman simplex and positivity & proved exact\\
Uniform fixed power-scale asymptotic & proved\\
State-wise RBLPTE in power sectors & proved\\
One-dimensional scalar transport & proved\\
Full target-preserving $\CSHT$ & not proved\\
Subpower core $\SPCC$ & open / RH-bearing\\
Root $\RBLPTE$ & open / RH-bearing\\
Riemann Hypothesis & unproved\\
\bottomrule
\end{tabular}
\end{center}

\end{document}
